\documentclass[../../main.tex]{subfiles}
\graphicspath{{../../image/}{../../figure/}}
\begin{document}

\section{强激波、弱激波与中等强度激波}\label{section:ch08-强激波弱激波中等强度激波}

对一维双曲系统
\begin{align}
\partial_t U+\partial_x F(U)=0,
\label{align:ch08-守恒律一维}
\end{align}
其 Rankine--Hugoniot 跳跃条件~\eqref{align:ch08-rh} 蜕化为
\begin{align}
F(U_+)-F(U_-)=s(U_+-U_-).
\label{align:ch08-rh}
\end{align}

实际上，一旦激波的传播方向 $v$ 被取定，且 $F$ 已按\eqref{align:ch07-平面前F}定义，\eqref{align:ch08-rh} 与多维版本\eqref{align:ch03-RH一般形式} 同样一般。

当\eqref{align:ch08-rh} 成立时，我们称状态 $U_-$（在左侧）以速度 $s$ 的激波与状态 $U_+$（在右侧）相连。注意，借助变换 $(x,t)\mapsto(-x,-t)$ 下\eqref{align:ch08-守恒律一维} 的不变性，\eqref{align:ch08-rh} 中“左”与“右”可以互换。尽管如此，后文我们将引入诱导不可逆性的容许性条件，其结果是 $U_-$ 与 $U_+$ 的角色不能再互换。

跳跃 $U_+-U_-$ 称为振幅，其大小 $|U_+-U_-|$ 称为激波的强度。不对激波强度施加限制而建立的性质，称为对强激波亦成立。然而，我们常常必须对激波强度施加限制：$|U_+-U_-|<\delta$，其中 $\delta$ 依赖于 $\mathrm{D}F$，经由诸如不同特征族的特征速度间距（它诱导不同族波的分离）以及函数 $\lambda_i$ 与 $R_i$ 的导数大小（它们刻画系统的非线性）等参数。特别地，当 $\delta$ 仅依赖于 $\lambda_i$ 与 $R_i$ 的一阶导数大小时，称激波为中等强度；而当 $\delta$ 还依赖于二阶导数大小时，称激波为弱激波。当然，这些参数的大小可通过 rescale 变量 $x,t$ 与 $U$ 而改变，因此真正相关的是 $\delta$ 的相对大小而非绝对大小。

注意到\eqref{align:ch08-rh} 可写成
\begin{align}
\big[A(U_-,U_+)-sI\big](U_+-U_-)=0,
\label{align:ch08-rh平均}
\end{align}
其中我们使用记号
\begin{align}
A(V,U)=\int_0^1 \mathrm{D}F\big(\tau U+(1-\tau)V\big)\,\mathrm{d}\tau.
\label{align:ch08-平均jacobi}
\end{align}

于是 $s$ 必须是 $A(U_-,U_+)$ 的实特征值，相应特征向量为 $U_+-U_-$。若对某个 $\bar U\in\mathcal O$，特征速度 $\lambda_i(\bar U)$ 是 $\mathrm{D}F(\bar U)$ 的单特征值，则当 $V,U$ 邻近 $\bar U$ 时，$A(V,U)$ 将有单实特征值 $\mu_i(V,U)$，相应特征向量为 $S_i(V,U)$。特别地，$A(U,U)=\mathrm{D}F(U)$，由此得 $\mu_i(U,U)=\lambda_i(U)$，$S_i(U,U)=R_i(U)$。注意 $A(V,U)$，从而 $\mu_i(V,U)$ 与 $S_i(V,U)$，关于 $(V,U)$ 对称。因此，将这些函数关于中点 $\frac12(V+U)$ 作（有限）Taylor 展开，得
\begin{align}
\mu_i(V,U)&=\lambda_i\!\left(\tfrac12(V+U)\right)+O(|V-U|^2),
\label{align:ch08-mu展开}\\
S_i(V,U)&=R_i\!\left(\tfrac12(V+U)\right)+O(|V-U|^2).
\label{align:ch08-S展开}
\end{align}

现设对一个中等强度激波
\begin{align}
s&=\mu_i(U_-,U_+),
\label{align:ch08-s中等}\\
U_+-U_-&=\zeta S_i(U_-,U_+).
\label{align:ch08-跳跃方向}
\end{align}
于是 $s$ 将接近特征速度 $\lambda_i$。这样的激波称为 $i$-激波。

\eqref{align:ch08-mu展开}、\eqref{align:ch08-s中等} 的一个有趣推论是有用的恒等式
\begin{align}
s=\tfrac12\big[\lambda_i(U_-)+\lambda_i(U_+)\big]+O\big(|U_--U_+|^2\big).
\label{align:ch08-s平均}
\end{align}

在特殊系统中，即使强激波也可关联到某个特定特征族。例如，等熵弹性系统\eqref{align:ch07-p系统} 的 Rankine--Hugoniot 条件
\begin{align}
\begin{cases}
v_+-v_-+s(u_+-u_-)=0\\
\sigma(u_+)-\sigma(u_-)+s(v_+-v_-)=0
\end{cases}
\label{align:ch08-弹性rh}
\end{align}

蕴涵
\begin{align}
s=\pm\sqrt{\frac{\sigma(u_+)-\sigma(u_-)}{u_+-u_-}}.
\label{align:ch08-弹性s}
\end{align}
回顾该系统的特征速度\eqref{align:ch07-p系统特征值}，自然地称向左传播（$s<0$）的激波为 1-激波，向右传播（$s>0$）的激波为 2-激波。

另一个重要例子是控制无黏性气体直线性绝热流的系统\eqref{align:ch07-热弹性}。其 Rankine--Hugoniot 跳跃条件为\footnote{此处 $r$ 表示激波速度，因为符号 $s$ 被保留用于表示比熵。}
\begin{align}
\begin{cases}
v_+-v_-+r(u_+-u_-)=0\\
p(u_+,s_+)-p(u_-,s_-)-r(v_+-v_-)=0\\
v_+p(u_+,s_+)-v_-p(u_-,s_-)-r\Big[\varepsilon(u_+,s_+)+\tfrac12 v_+^2-\varepsilon(u_-,s_-)-\tfrac12 v_-^2\Big]=0.
\end{cases}
\label{align:ch08-气体rh}
\end{align}
关联于特征速度 $\lambda_2=0$ 的 2-激波是定常的：$r=0$，此时\eqref{align:ch08-气体rh} 蜕化为 $v_-=v_+$ 与 $p_-=p_+$。另一方面，1-激波与 3-激波以负速度和正速度传播：
\begin{align}
r=\pm\sqrt{-\frac{p(u_+,s_+)-p(u_-,s_-)}{u_+-u_-}}.
\label{align:ch08-气体r}
\end{align}

进而，当 $r\neq0$ 时，可将\eqref{align:ch08-气体rh} 中三个方程联立，推导出著名的 Hugoniot 方程
\begin{align}
\varepsilon(u_+,s_+)-\varepsilon(u_-,s_-)=-\tfrac12\big[p(u_+,s_+)+p(u_-,s_-)\big](u_+-u_-)
\label{align:ch08-hugoniot方程}
\end{align}
（见历史引论），它不含速度或激波速度，而仅关联热力学状态变量 $u$ 与 $s$。

\end{document}
